\documentclass[12pt,a4paper]{article}
\usepackage[utf8]{inputenc}
\usepackage[T1]{fontenc}
\usepackage{lmodern}
\IfFileExists{brazil.ldf}{\usepackage[brazil]{babel}}{%
  \usepackage{babel}\babelprovide[main,import]{portuguese}}
\usepackage[margin=2.5cm,headheight=16.6pt]{geometry}
\usepackage{array,tabularx,booktabs}
\usepackage{amsmath}
\usepackage{enumitem}
\usepackage{fancyhdr}
\usepackage{xcolor}
\usepackage{titlesec}
\definecolor{azulProfessor}{HTML}{184C78}
\definecolor{cinzaTitulos}{HTML}{1E1E1E}
\definecolor{cinzaCabecalho}{gray}{0.5}
\titleformat{\section}{\Large\bfseries\color{azulProfessor}}{\thesection.}{0.6em}{}[\titlerule]
\titleformat{\subsection}{\large\bfseries\raggedright\color{cinzaTitulos}}{\thesubsection.}{0.6em}{}
\titleformat{\subsubsection}{\normalsize\bfseries\raggedright\color{cinzaTitulos}}{\thesubsubsection.}{0.6em}{}
\titlespacing*{\section}{0pt}{3.5ex plus 1ex minus .2ex}{2.3ex plus .2ex}
\titlespacing*{\subsection}{0pt}{3.25ex plus 1ex minus .2ex}{1.5ex plus .2ex}
\titlespacing*{\subsubsection}{0pt}{2.5ex plus 1ex minus .2ex}{1.2ex plus .2ex}
\setcounter{tocdepth}{2}
\usepackage{hyperref}
\hypersetup{colorlinks=true,linkcolor=azulProfessor,urlcolor=azulProfessor,citecolor=azulProfessor,linktoc=section,pdftitle={Laboratório Educativo de IA Generativa com GANs — Relatório Final — Sprint 3},pdfauthor={Allan Vitor Piran Lemes e Guilherme Rodrigues dos Santos}}
\urlstyle{same}
\newcolumntype{Y}{>{\raggedright\arraybackslash}X}
\setlength{\parindent}{1.5em}
\setlength{\parskip}{0pt}
\setlength{\emergencystretch}{2em}
\renewcommand{\arraystretch}{1.2}
\setlist[itemize]{leftmargin=*,itemsep=5pt,topsep=6pt}
\pagestyle{fancy}
\fancyhf{}
\fancyhead[L]{\small\color{cinzaCabecalho}Oficina de Projetos em Ciência de Dados e IA}
\fancyhead[R]{\small\color{cinzaCabecalho}Relatório Final -- UTFPR}
\fancyfoot[L]{\small\color{cinzaCabecalho}Documento de Ciência de Dados \& IA}
\fancyfoot[R]{\small Página \thepage}
\renewcommand{\headrulewidth}{0.4pt}
\renewcommand{\footrulewidth}{0.4pt}
\newcommand{\pendente}[1]{\textit{#1}}

\begin{document}
\renewcommand{\contentsname}{Sumário}
\renewcommand{\refname}{Referências}

\begin{titlepage}
\thispagestyle{fancy}
\centering
\vspace*{0.7cm}
{\Large\bfseries UNIVERSIDADE TECNOLÓGICA FEDERAL DO PARANÁ\par}
\vspace{0.15cm}
{\large DEPARTAMENTO ACADÊMICO DE INFORMÁTICA\par}
\vspace{0.1cm}
{\small CAMPUS PATO BRANCO\par}
\vspace{0.1cm}
{\small OFICINA DE PROJETOS EM CIÊNCIA DE DADOS E IA\par}
\vspace{1.6cm}
\noindent\rule{\textwidth}{1pt}\par
\vspace{0.2cm}
{\Huge\bfseries\color{azulProfessor}Laboratório Educativo de IA\par Generativa com GANs\par}
\vspace{0.35cm}
{\large\bfseries RELATÓRIO FINAL DE PROJETO\par}
\vspace{0.25cm}
{\small Versão de trabalho: Sprint 3 e ajuste de escopo\par}
\vspace{0.45cm}
\noindent\rule{\textwidth}{1pt}\par
\vspace{1.2cm}
{\renewcommand{\arraystretch}{1.05}
\begin{tabular}{@{}ll@{}}
\textbf{Equipe de Desenvolvimento:} & Allan Vitor Piran Lemes \\
 & Guilherme Rodrigues dos Santos \\
\textbf{Orientador / Prof.:} & Dr. Luiz Fernando Puttow Southier \\
\textbf{Data de Atualização:} & 15 de novembro de 2026 \\
\textbf{Versão do Documento:} & 3.0 \\
\end{tabular}\par}
\end{titlepage}
\setcounter{page}{2}
\tableofcontents
\clearpage

\section{Introdução e ajuste de escopo}
O Laboratório Educativo de IA Generativa com GANs é desenvolvido por Allan Vitor Piran Lemes e Guilherme Rodrigues dos Santos na disciplina Oficina de Projetos em Ciência de Dados e IA, da UTFPR. O projeto reúne uma GAN com Fashion-MNIST, aplicação local e materiais introdutórios, com linguagem acessível ao ensino médio e a leitores iniciantes.

\textbf{Em 15/11/2026, a equipe solicitou o ajuste da entrega educativa para produção de materiais e verificação pelos próprios autores.}

O relatório anterior está em \path{docs/Sprint2.tex}. Esta versão mantém o histórico técnico, atualiza a forma de entrega e não substitui fatos anteriores por atividades fictícias. Os arquivos Markdown das sprints são memória de trabalho entre chats; o presente fonte LaTeX é a versão corrente de trabalho para entrega acadêmica.

\section{Objetivos e entregáveis}
O objetivo geral é desenvolver material educativo introdutório sobre IA generativa utilizando uma GAN treinada com Fashion-MNIST e uma aplicação local de demonstração. A entrega inclui apostila com explicações acessíveis, exemplos visuais, orientações de uso e aprendizados documentados da equipe.
\begin{itemize}
\item Compreender o papel do gerador, do discriminador e dos dados.
\item Preparar os dados e adaptar uma implementação introdutória.
\item Treinar a GAN, salvar o modelo e registrar a evolução visual.
\item Disponibilizar aplicação local para conhecer a base, comparar épocas e gerar imagens.
\item Produzir apostila, roteiro de exploração autônoma e documentação acadêmica.
\item Verificar funcionamento e coerência dos materiais, com revisão pelos autores.
\end{itemize}
A interface não oferece escolha de classe nem download. A exportação de PNG é realizada por um script de terminal. A arquitetura incondicional recebe ruído, sem texto ou rótulo de categoria.

\section{Metodologia e critérios de verificação}
O trabalho combina estudo introdutório, implementação incremental, testes técnicos e análise visual descritiva. O Fashion-MNIST contém imagens de roupas e acessórios de $28\times28$ pixels em tons de cinza~\cite{fashionmnist}. A preparação separa as 60.000 imagens oficiais de treino em 50.000 para treino e 10.000 para validação, com estratificação e semente 42. As 10.000 imagens oficiais de teste permanecem separadas. Os pixels são normalizados para $[-1,1]$; validação e teste não atualizam as redes.

O gerador transforma um vetor aleatório de dimensão 100 em imagem. O discriminador aprende a distinguir imagens reais das geradas. O treinamento alterna atualizações com Adam e perda binária. As perdas não são porcentagens de acerto. As imagens com ruído fixo permitem observar a evolução dentro de uma execução. O checkpoint guarda redes, otimizadores e estados aleatórios ao final da época; o gerador também é exportado separadamente.

A conclusão técnica depende de arquivos existentes, execução dos comandos e correspondência entre exemplos e explicações. A revisão autoral deve conferir linguagem, limitações, identificação das imagens, instruções de uso e coerência do relatório. Os testes automáticos não substituem essa revisão; sua conclusão não foi informada pela equipe.

\section{Histórico técnico das sprints}
\subsection{Sprint 1: organização e preparação}
O registro anterior descreve a organização do repositório e do quadro inicial de tarefas, além do desenvolvimento da preparação dos dados. Foram documentadas a divisão estratificada, a distribuição de 5.000 imagens por classe no treino e 1.000 na validação e a exportação de exemplos reais. As contribuições informadas foram organização do repositório, dados e caracterização do destinatário por Allan; organização do quadro e planejamento por Guilherme.

O relatório anterior informa 3 horas de Allan e 2 de Guilherme, totalizando 5 horas diante de 10 previstas. Esses valores são preservados como informações anteriores da equipe, sem nova apuração nesta revisão. O desvio informado foi de $-50\%$, sem que isso comprove produtividade ou conclusão de todos os itens.

\subsection{Sprint 2: implementação e primeiro treinamento}
O commit \texttt{0523e6c}, de Allan, contém a implementação inicial e grades identificadas como épocas 0, 1, 5 e 15. O treinamento original carregava 60.000 imagens. Na adaptação local, registrada em \texttt{451cde6}, a arquitetura foi mantida e o treino passou a usar a divisão de 50.000 imagens. Foram acrescentados retomada, registros, visualização em modo de avaliação e geração a partir do modelo salvo.

\begin{center}
\begin{tabularx}{\textwidth}{@{}YY@{}}
\toprule
\textbf{Execução local documentada} & \textbf{Resultado} \\
\midrule
Identificação & \path{sprint2_local} \\
Épocas completas / imagens de treino & 1 / 50.000 \\
Lote máximo / semente & 128 / 42 \\
Adam: taxa / coeficientes & 0,0002 / 0,5 e 0,999 \\
Processamento & CPU, duas threads \\
Tempo do laço & 296,8 segundos \\
Perdas médias D / G & 0,7726 / 1,5397 \\
\bottomrule
\end{tabularx}
\end{center}

Esses resultados pertencem à Sprint 2 e continuam sendo os últimos comprovados pelo histórico local. O tempo não inclui download e exportação, nem equivale a horas ativas dos autores. As grades antigas de 15 épocas são de outra execução, sem checkpoint correspondente comprovado neste checkout.

Foram verificados separação dos dados, carregamento do gerador, retomada equivalente ao treino contínuo em CPU e preservação das estatísticas da rede durante a visualização. As imagens locais iniciais ainda apresentam borramento e padrões quadriculados. Não há métrica quantitativa de qualidade implementada.

\subsection{Sprint 3: aplicação e materiais}
O registro fornecido pela equipe informa seis itens concluídos no Jira, com prazo planejado de 13/11/2026. A revisão do checkout em 15/11/2026 encontrou o modelo local da época 1 e ausência de interface e materiais específicos da Sprint 3. Esses artefatos foram implementados durante a intervenção registrada, sem atribuir automaticamente aos responsáveis informados a autoria de cada alteração. O Jira não foi acessado nem alterado neste ajuste.

\begin{center}
\small
\begin{tabularx}{\textwidth}{@{}lYY@{}}
\toprule
\textbf{Item} & \textbf{Responsável informado} & \textbf{Evidência atual} \\
\midrule
SCRUM-11 & Guilherme & Gerador utilizável da época 1; continuidade do treino não comprovada. \\
SCRUM-12 & Allan & Comparação de épocas originais 0, 5 e 15; execução local separada. \\
SCRUM-13 & Guilherme & Aplicação Streamlit com exemplos, comparação e geração. \\
SCRUM-14 & Allan & Explicações sobre dados, redes e limitações; revisão autoral final a registrar. \\
SCRUM-26 & Guilherme & Fluxograma técnico e diagrama simplificado na aplicação. \\
SCRUM-15 & Allan & Roteiro autônomo, memória técnica e organização documental. \\
\bottomrule
\end{tabularx}
\end{center}

As estimativas informadas para SCRUM-11, 12, 13, 14 e 15 são, respectivamente, 5, 4, 5, 3 e 3 horas. Somam 20 horas previstas, com 10 para cada integrante; SCRUM-26 não possui estimativa informada. As horas efetivas das Sprints 2 e 3 e reflexões individuais não foram fornecidas. Não se calcula desvio dessas etapas sem esses dados.

A aplicação permite mudar a semente e gerar 16 imagens com o modelo salvo. Os testes funcionais verificaram abertura, troca da comparação, geração, formato, intervalo dos pixels, repetição por semente, alteração por nova semente e mensagem de modelo ausente. O servidor respondeu à verificação de funcionamento. Não houve inspeção visual em navegador.

\section{Materiais educativos e uso}
A apostila \path{docs/apostila_educativa.tex} apresenta conceitos, imagens identificadas, orientações e questões comentadas. O roteiro \path{docs/roteiro_sprint3.md} orienta leitura autônoma por aproximadamente 25 minutos e pode apoiar a apresentação acadêmica ao professor. O fluxograma está em \path{docs/fluxograma.md}; a aplicação contém uma versão simplificada.

Na raiz do repositório, após instalar as dependências:
\begin{quote}
\small\verb|python -m streamlit run app.py --server.address 127.0.0.1|
\end{quote}
O endereço local é \url{http://localhost:8501}. A geração exige \path{checkpoints/sprint2_local/gerador.pt}, fora do Git. Os exemplos salvos podem ser consultados sem executar outro treinamento. As grades originais usam contraste automático e modo de treino; as locais usam escala fixa e modo de avaliação. A aplicação mantém as execuções separadas.

\section{Riscos, limitações e próximos passos}
\begin{itemize}
\item Dificuldade conceitual: manter explicações introdutórias, conferir o código e revisar a apostila pelos autores.
\item Treino demorado ou interrompido: testar antes, salvar por época e retomar o checkpoint. Uma época incompleta será refeita.
\item Imagens ambíguas: apresentar limitações e observar evolução sem afirmar qualidade ou acurácia inexistentes.
\item Falta de tempo: priorizar artefatos executáveis, revisão dos materiais e consolidação do relatório.
\item Divergência documental: distinguir execução original e local, estimativas e horas efetivas, estado informado no Jira e evidência disponível.
\end{itemize}
Permanecem necessários a confirmação de eventual modelo posterior ou o registro da continuação do treinamento, a revisão final dos materiais pelos autores, as horas e reflexões individuais e a consolidação do relatório final. Não foi executado treinamento adicional para este ajuste de escopo. A geração por classe não integra a implementação atual.

\section{Conclusão parcial}
O projeto dispõe de uma GAN executável, checkpoint local de uma época, registros visuais, aplicação Streamlit e materiais introdutórios. A entrega educativa passou a se concentrar na apostila, documentação e exploração local. A mudança altera a forma de entrega, sem ampliar artificialmente os resultados técnicos. O fechamento acadêmico deve refletir as evidências e a revisão dos autores, preservando o histórico e as limitações descritas.

\begin{thebibliography}{9}
\addcontentsline{toc}{section}{Referências}
\bibitem{fashionmnist} XIAO, Han; RASUL, Kashif; VOLLGRAF, Roland. \textbf{Fashion-MNIST: a Novel Image Dataset for Benchmarking Machine Learning Algorithms}. 2017. \url{https://github.com/zalandoresearch/fashion-mnist}.
\bibitem{pytorch} PYTORCH. \textbf{DCGAN Tutorial}. \url{https://docs.pytorch.org/tutorials/beginner/dcgan_faces_tutorial.html}.
\bibitem{checkpoint} PYTORCH. \textbf{Saving and Loading Models}. \url{https://docs.pytorch.org/tutorials/beginner/saving_loading_models.html}.
\bibitem{registros} EQUIPE DO PROJETO. Registros técnicos internos das Sprints 2 e 3 e artefatos de execução. Arquivos locais de apoio; não substituem os relatórios oficiais.
\end{thebibliography}
\end{document}
